\chapter{The Prior Art: What Is Taken and What Is Declined}
\label{app:priorart}

Chapter~\ref{sec:related} compares the design with related work. This appendix records, line of
work by line of work, exactly what the design takes from each, what it declines, and why.

\section{Clean}

\emph{Sources:} Barendsen and Smetsers \citep{barendsen1993conventional,barendsen1995inference,barendsen1996uniqueness},
Smetsers et al.\ \citep{smetsers1994guaranteeing}, de Vries et al.\
\citep{devries2007redefined,devries2008simplified}, and the Clean language report
\citep{plasmeijer2011clean}.

\emph{Taken.} The demand lives in library types: \code{newArray} returns an array whose uniqueness
attribute is a variable rather than ``unique'' \citep[\S6]{devries2008simplified}, so user code carries
nothing. The propagation rule for containers---``if a unique object is stored in a data structure, the
data structure itself becomes unique as well'' \citep[\S9.2]{plasmeijer2011clean}---appears here as the
holder relation through paths. Evaluation-order-aware marking, Smetsers et al.'s ``preliminaries'' and
groups of ``alternate arguments'' \citep[\S4]{smetsers1994guaranteeing} and the report's observing references
\citep[\S9.4]{plasmeijer2011clean}, is what liveness at the site generalises (\S\ref{sec:model-unique}).
Deep uniqueness is the source of the exclusive-contents flag (\S\ref{sec:model-excl}).

\emph{Declined.} Attributes on types solved by unification with a subtyping restriction on arrows
\citep[Def.~7.6]{barendsen1996uniqueness}, for which Barendsen and Smetsers give no principal
uniqueness type theorem \citep[\S8]{barendsen1996uniqueness}; the rule that the type of every binding in
a recursive binding group is non-unique \citep[\S7]{devries2008simplified}, which concerns the
recursively bound names rather than their parameters; and the rank-2 escape that de Vries et al.\ use to
type a function such as \code{if isEmpty arr then shrink arr else grow arr}
\citep[\S6]{devries2008simplified}, which an annotation-free checker cannot use and liveness makes
unnecessary. One of Clean's known false errors, a graph-marking function whose array lookup precedes its
update and which is rejected because the array ``is used twice in the right-hand side''
\citep[\S6]{barendsen1995inference}, is accepted here, because the lookup is complete before the update;
Clean's remedy, a \code{lookup} that returns the array, is unnecessary.

\section{Usage aspects}

\emph{Sources:} Aspinall and Hofmann \citep{aspinall2002another}; Aspinall, Hofmann and
Kone\v{c}n\'y \citep{aspinall2008usage}; Kone\v{c}n\'y \citep{konecny2003typing}, through Aspinall et
al.'s report of it.

\emph{Taken.} The three aspects \citep[\S1]{aspinall2008usage} as the \textsf{W}, \textsf{A} and
\textsf{R} facts of $\mathit{use}$; the \textsc{Let} rule's per-variable side condition, under which a
variable used with aspect 3, ``used read-only and not shared with the result'', in the binding may be
destroyed in the body, while one used with aspect 2 may not \citep[Fig.~8]{aspinall2008usage}---the
reads-before-writes rule in its earliest typed form, restricted to reads whose result does not alias the
variable, which is also why, in our rules, a read whose result aliases the list keeps a holder alive; the
best-typing claim and its optimistic search \citep[\S6.2]{aspinall2008usage}; the
separation theorem's clause that heap occupied by arguments with aspect 2 or 3 is not modified
\citep[Thm.~5.4]{aspinall2008usage}, as clauses 1 and 2 of Lemma~\ref{lem:summary}; and the
two-semantics soundness theorem \citep[Thm.~5.8]{aspinall2008usage} as the shape of
Theorem~\ref{thm:main}.

\emph{Declined.} The system is first-order, monomorphic and affine, with explicit resources for
allocation; higher-order functions are deferred there to the standard types-and-effects technique,
which the ownership classes of \S\ref{sec:inf-classes} are.

\section{Linearity, alias burying and the past}

\emph{Sources:} Wadler \citep{wadler1990linear,wadler1991use}; Boyland \citep{boyland2001alias};
Marshall, Vollmer and Orchard \citep{marshall2022entente}.

\emph{Taken.} The hyperstrict ordering of \code{let!} as the reason strictness matters
\citep{wadler1990linear}; ``only update operations are sequentialised but lookups can occur in
parallel'' \citep[\S3]{wadler1991use} as the design goal; alias burying's requirement that all aliases
be dead \citep[\S3.2]{boyland2001alias} as the definition of unique at a site; and uniqueness as ``a
guarantee about what has been done with a value `in the past'\,'' \citep[\S2.2]{marshall2022entente}. In
their calculus ``there is no way to transform a non-unique value into a unique one'' (ibid.); regaining
uniqueness after a borrow appears there only as future work, by recombining fractional permissions
\citep[\S6]{marshall2022entente}. Our checker gets the same effect without fractions: uniqueness is
decided at a program point, so a value whose other holders are all dead is unique again at a later site,
and \code{List.copy} is the way back only when a holder stays live.

\emph{Declined.} Boyland's annotations on fields and interfaces, which are inferred here; Marshall et
al.'s call-by-name calculus.

\section{Destructive-update analysis}

\emph{Sources:} Wand and Clinger \citep{wand2001set}; the Sisal compiler \citep{cann1990sisal}.

\emph{Taken.} Liveness-defined destructive update---an update is destructive when the array ``can't be
bound to a live location'' after its arguments have been evaluated---and inclusion constraints with a
least solution, which ``can be solved in polynomial time'' \citep{wand2001set}; Sisal's tabulation of
copy sites as the model for the hand classification of \S\ref{sec:eval-hand}.

\emph{Declined.} Wand and Clinger's first-order restriction to arrays of scalars, and both systems'
fallback to a copy.

\section{Reference counting with reuse}

\emph{Sources:} Lean~4 \citep{ullrich2019beans,demoura2021lean4}, Perceus \citep{reinking2021perceus},
frame-limited reuse \citep{lorenzen2022frame}, FP\textsuperscript{2} \citep{lorenzen2023fp2},
Lorenzen et al.\ on functional in-place algorithms \citep{lorenzen2024bst}, and Brandon et al.\
\citep{brandon2026borrows}.

\emph{Taken.} The direction of inference, from ``all borrowed'' towards a fixed point
\citep[\S5.2]{ullrich2019beans}, and Brandon et al.'s least fixpoint justified by Knaster--Tarski
\citep[\S4.4]{brandon2026borrows}; the one correctness rule every paper shares, that a parameter that
reaches a write must be owned; Brandon et al.'s refinement that treats variable occurrences in the
argument position of a tail call as escaping, which beni's loop rebinding handles directly; and
FP\textsuperscript{2}'s diagnosis of static uniqueness typing as ``code duplication, where a single
function can have multiple different implementations'' \citep[\S1.4.1]{lorenzen2023fp2}, answered here
by summaries that are facts about flows, not types, so one \code{reverse} serves both kinds of caller
and the \emph{caller} is what is checked.

\emph{Declined.} Run-time reference counts and the silent copy at a shared call---``the fip annotation
in Koka only guarantees that no (de)allocation occurs if the parameters are unique at runtime''
\citep[\S3.1]{lorenzen2024bst}; borrowing as a discipline on reference-count tokens that lets a borrowed
value be returned with an increment \citep{ullrich2019beans}. Lorenzen and Leijen's objection that
inferred borrowing can extend a value's lifetime and so cost space \citep{lorenzen2022frame} does not
apply, because there is no count to keep a borrowed value alive past its last read. Brandon et al.'s
fallback, inserting ``a handful of reference count operations'' where a program cannot be typed, is
this design's error.

\section{OxCaml}

\emph{Sources:} Lorenzen et al.\ \citep{lorenzen2024oxidizing}, Peters et al.\
\citep{peters2026mode}, and the OxCaml documentation \citep{oxcaml2025uniqueness}.

\emph{Taken.} The lock rule for closures \citep[\S3.3]{lorenzen2024oxidizing} as the $\mathit{once}$
rule of \S\ref{sec:hard-closures}; polarised inequality constraints solved by what is essentially a
transitive closure as the shape of the edges of \S\ref{sec:inf-classes}; mode crossing
\citep{peters2026mode} as $\mathit{crosses}(\tau)$; the measured count of 27\,382 locality annotations in
interfaces, where ``the implementations can generally infer'' \citep[\S7]{lorenzen2024oxidizing}, as the
warning about interface churn; and the remark that the complications of currying go away in a language
without it \citep[\S6.4]{lorenzen2024oxidizing} as evidence that beni's decision removes the silent
capture through partial application---explicit closures still need the lock rule, and the once closures of
\S\ref{sec:model-functions} are that rule.

\emph{Declined.} Modes as types with no polymorphism over them, which led to functions duplicated per
mode \citep[\S6.6]{lorenzen2024oxidizing}---summaries here are per flow, not per mode; and per-module
inference with modes written in interfaces---summaries here are inferred \emph{and} published.

\section{Reachability types and capture tracking}

\emph{Sources:} Bao et al.\ \citep{bao2021reachability}, Wei et al.\ \citep{wei2024polymorphic}, Jia et
al.\ \citep{jia2026escape}, O'Connor et al.\ \citep{oconnor2026separation}, System Capybara
\citep{xu2026capybara}, and Deng et al.\ \citep{deng2025capture}.

\emph{Taken.} A function type that names which argument its result reaches \citep[\S2.2]{wei2024polymorphic}
as $\mathit{res}$; the warning that polymorphism over tracked and untracked values is unsound without a
freshness marker---here $\fresh$ is a distinct kind of cell, never a type variable; Capybara's rule that
consuming a value kills its aliases as the content of the dead-alias lemma, and its rule that a component
of a shared structure cannot be consumed directly as the containers rule of \S\ref{sec:hard-containers};
Deng et al.'s use and kill effects on closures \citep[Fig.~1]{deng2025capture} as the cells a closure
holds under $\env$ and the once-ness of a closure that writes them, inferred rather than declared; and
O'Connor et al.'s statement of the guarantee.

\emph{Declined.} Qualifier annotations on function arguments and explicit instantiation, which
bidirectional inference for reachability types still requires \citep{jia2026escape}; and an opaque joint
qualifier for the elements of a list, replaced here by per-path cells.

\section{Rust}

\emph{Sources:} non-lexical lifetimes \citep{rfc2094}, two-phase borrows \citep{rfc2025}, Polonius
\citep{matsakis2018alias,stjerna2020polonius}, Oxide \citep{weiss2019oxide}, Pearce's lightweight formalism
\citep{pearce2021lightweight}, Ho et al.'s symbolic semantics \citep{ho2024sound}, and Emre et al.\
\citep{emre2023aliasing}.

\emph{Taken.} Liveness as the precision lever, and the minimal-lifetime principle \citep{rfc2094};
two-phase borrows as the formal statement of reads before writes \citep{rfc2025}; the three-point error
with the blame on the action \citep{rfc2094}; Polonius as transitive closures over directional facts
\citep{matsakis2018alias,stjerna2020polonius}; Pearce's and Ho et al.'s treatment of joins at merge points as the model for
\S\ref{sec:hard-branches}; and Emre et al.'s measurement as the argument against unification
\citep{emre2023aliasing}.

\emph{Declined.} Lifetimes in signatures, and references into structures; beni has none, which is why
the third problem case of non-lexical lifetimes does not arise.

\section{Futhark}

\emph{Sources:} Henriksen et al.\ \citep{henriksen2017futhark,henriksen2017thesis} and the Futhark
blog and documentation \citep{henriksen2021anatomy,henriksen2022uniqueness,henriksen2026aliasing,henriksen2026rewrite,futhark2026errorindex}.

\emph{Taken.} The consumption judgement $(\mathcal{O}_2 \cup \mathcal{C}_2) \cap \mathcal{C}_1 = \emptyset$
\citep{henriksen2017futhark} as the one-line statement of the rule; the union of alias sets for a
conditional, kept for cells, with per-path liveness added; the refusal of a consumption of a free
variable inside a lambda \citep[\S3.3]{henriksen2017futhark} as the $\mathit{once}$ rule; the separate
pass after type checking \citep{henriksen2026rewrite}; the error index's shape---message, minimal program,
reason, a copy as the fix, and a note when the copy ``subverts the purpose''
\citep{futhark2026errorindex}; the lesson on naming intermediate results \citep{henriksen2021anatomy};
and the free-theorem idea for the identity function and the pipe \citep{henriksen2026aliasing}.

\emph{Declined.} Uniqueness annotations on parameters and results; an intraprocedural analysis with no
summaries, for which ``we are forced to be conservative'' \citep[\S3.2]{henriksen2017futhark}; and
Henriksen's 2026 advice that other designers avoid the feature in favour of reference counting with
copy-when-shared \citep{henriksen2026aliasing}. We decline that advice for the reason he gives for
Futhark itself: a guarantee is wanted, here that a written in-place edit does not copy. The objections he
raises---abstraction hiding aliasing, interface clutter, spurious errors through generic helpers---apply
to inferred summaries as well; \S\ref{sec:model-cells} (opaque types), \S\ref{sec:inf-modular} (churn),
\S\ref{sec:hard-pipelines} (helpers) and \S\ref{sec:disc-runtime} (the run-time alternative) answer
them as far as a design can.

\section{Usage analysis, affine closures and escape analysis}

\emph{Sources:} Turner et al.\ \citep{turner1995once}, Wansbrough and Peyton Jones
\citep{wansbrough1999once,wansbrough2002thesis}, Affe \citep{radanne2020kindly}, and Blanchet
\citep{blanchet1998escape}.

\emph{Taken.} The poisoning problem \citep[\S6.3]{wansbrough1999once} as a failure mode summaries must
avoid (they solve it by subsumption, and argue against usage polymorphism); their observation that only
usage polymorphism expresses dependencies among arguments, as in \code{apply}, as the reason our summaries
are polymorphic over function-typed parameters;
polarity-based simplification of constraint sets \citep{wansbrough2002thesis} as the way to keep summary
blocks small; the measured warning that generalising local binders loses precision; Affe's rule under
which a closure that captures a shared borrow stays unrestricted \citep[\S3.3]{radanne2020kindly} as the
shape of a closure that only reads its environment; and Blanchet's case split for closures, applied
versus escaped \citep{blanchet1998escape}, as the $\mathsf{E}$ fact on a function-typed parameter's
environment.

\emph{Declined.} Usage on every type, which counts a read-only capture as a use; Affe's lexical regions
and explicit borrows.

\section{Error design}

\emph{Sources:} Czaplicki \citep{czaplicki2015errors}, Wrenn and Krishnamurthi \citep{wrenn2017classifiers},
Marceau et al.\ \citep{marceau2011mind}, the rustc development guide \citep{rustc2026diagnostics},
Crichton et al.\ \citep{crichton2020usability,crichton2023grounded,crichton2024profiling}, Hylo
\citep{hylo2024bindings}, and Swift \citep{swift2025explicitcopies,swift2026diagnosticssil}.

\emph{Taken} in full in Chapter~\ref{sec:errors}. The one dissent recorded is Marceau et al.'s advice
that error messages for novices should not propose solutions, against which Hylo, Futhark, Rust and
Swift all ship fixes; this design ships the fix and the verdict line (\S\ref{sec:errors-shape}).

\section{The prototype study}

The static variant of the prototype study (\S\ref{sec:bg-measure}; method in Appendix~\ref{app:method})
approximates this checker's backend half: it wrote in place where last use and per-function interface
summaries proved a value unique, with parameters consumed uniquely and the aliasing of results recorded in
interfaces and flowing forward to callers. It had no exclusive-contents flag, no ownership classes and no
entry protocol, so it is evidence for what in-place writes buy, not for these rules. Two of its findings
shaped the design: that the TEA model needs a runtime that hands the model over, and that a copy at a call
where uniqueness cannot be proved is catastrophic in a loop. They are why the rule is an error at the site
and not a copy. A third finding cuts the other way: a run-time shared bit made the same in-place decisions
(\S\ref{sec:disc-runtime}).

\section{Compile-time garbage collection for Mercury}

\emph{Sources:} Mazur \citep{mazur2004thesis} and Mazur et al.\ \citep{mazur2001practical}; earlier
sharing and update analyses known to us from secondary sources \citep{bloss1989update,hudak1987semantic,jones1989compile,baker1990unify}.

\emph{Taken.} Liveness at the site of reuse; path-based sharing facts with recursive types folded into
finite type graphs, as our folded paths; per-procedure summaries in interface files and a precondition
published for callers, as our summaries and side conditions.

\emph{Declined.} Commutative sharing pairs, replaced by directional cell sets; two compiled versions of
each procedure with a silent fallback to the unoptimised one, replaced by one body and an error; and the
exclusion of higher-order calls, replaced by ownership classes.
